\scanpage{15}
Il ne semble pas qu'on puisse obtenir des normalisations locales sans \deleted{\unsure{hyp}} restriction de car. qui fasse intervenir moins de paramètres, [si on veut qu'elle soit subordonnée à la descente \deleted{\unsure{si|sur}} par morphismes plats, ou même simples]. Si on admet des ext. de base finies et plates, on peut pousser la normalisation plus loin en exigeant que \deleted{\unsure{le}} $\delta(b,c,d,e,f)=1$ i.e. \inserted{(intrinsèquement)} se bornant à des trivialisations de $L$ \deleted{compatibles} définies par des sections $\lambda$ dont la puissance $12$.\textsuperscript{ème} (ou est-ce $24$.\textsuperscript{ème}) est \inserted{$\delta^{-1}$, où $\delta$ est} le discriminant. On a \deleted{\unsure{en car.}} alors des \deleted{\unsure{plus}} normalisations plus poussées en car $\ne2$, puis en car. $\ne3$, puis \unsure{on vieillira|on veillera}.
% archive sheet number: 14
